\chapter{M12: getrennte RANS-Code- und Modellprüfung}\label{vv:rans}
Der vorab festgelegte Vertrag aus Dokument 71 trennt Punkt-/MMS-Verifikation,
kontinuierliche Modellgleichungen und unabhängige DNS. Dokument 89 und das
versionierte M12-Nachweisarchiv nennen Quellenhashes und sämtliche Messwerte.
Die SA-MMS auf 16, 32, 64 und 128 Zellen erreicht relative Fehler
\(9.506\cdot10^{-4},2.482\cdot10^{-4},6.366\cdot10^{-5},1.613\cdot10^{-5}\)
und Ordnungen 1.937, 1.963 und 1.981. Der virtuelle Wandabstand dient allein diesem Test.

Ein unabhängiger kontinuierlicher Randwertlöser für den SA-Kanal liefert
die Codevergleichslösung. Die Netzfolge liefert relative RMS-U-Fehler:
\begin{center}
\begin{tabular}{rr}
\toprule
Zellen je Halbkanal & Fehler [Prozent] \\
\midrule
32 & 0.3473 \\
64 & 0.1874 \\
128 & 0.02123 \\
256 & 0.00882 \\
\bottomrule
\end{tabular}
\end{center}
Das feinste Gate von 0.5 Prozent ist erfüllt; Wandkraftfehler bleiben unter
\(1.8\cdot10^{-13}\), maximales \(y^+\) sinkt von 0.498 auf 0.0585.
Viskositätssensitivität, Wandverteilung und Relaxationswiederholung bestehen
bei unveränderten Modellkonstanten und absoluten Feld- und Bilanzgates.

Die unabhängige DNS von Moser/Kim/Mansour nutzt \(Re_\tau=587.19\),
nicht den gerundeten Dateinamen 590. Auf 128/256 Halbkanalzellen betragen
RMS-U+-Abweichungen 1.0775/1.0746 Prozent und Bulkabweichungen
0.3949/0.4074 Prozent; die Grenzen sind jeweils 5 Prozent.
Die letzte Profiländerung ist 0.01628 Prozent, kleiner als 1 Prozent und
ein Fünftel des Modellbudgets. Keine Anpassung von SA-Konstanten.

Das Quellenmanifest verlinkt Originaldaten und ihren SHA-256. Fehlende explizite
Weitergabelizenz verhindert ihre Redistribution; CI lädt sie nur temporär.
Die Profile sind DNS, keine experimentelle Messung. Dies ist eine begrenzte
Modellbeurteilung für den glatten voll entwickelten Kanal und keine allgemeine
RANS- oder FSI-Validierung. Ein anfänglicher Interpolationsfehler durch
Koordinatenrundung ist mit unabhängigem parabolischem Regressionstest korrigiert;
Solver, Originaldaten und Gates bleiben unverändert. Unabhängiges menschliches
Fachreview liegt nicht vor. M12 ist weiterhin ein Entwicklungskandidat.
